
@@M0@@
Comme @@M1@@, c'est la formule qui sert à définir @@M2@@.

ii) On se place dans le cas général. On compose alors @@M3@@ avec une section : @@M4@@, puis avec une section : @@M5@@, de façon que l'immersion composée vérifie la condition de 4.1. (voir 4.2). La proposition 4.3 est alors vraie pour l'immersion composée d'après ce qui précède; il reste donc à « redescendre » à @@M6@@. On redescend d'abord à @@M7@@, puis à @@M8@@, de sorte que la démonstration se fait en appliquant deux fois le lemme suivant (dû à J.P. Jouanolou) :

\textbf{Lemme} : \uline{Dans les conditions de} 4.3, \uline{soit} @@M9@@ \uline{la section définie en} 4.2. \uline{Si la proposition} 4.3 \uline{est vraie pour} @@M10@@, \uline{elle est vraie pour} @@M11@@.

Si on suppose 4.3 vraie pour @@M12@@, on a pour tout @@M13@@, et tout @@M14@@ :
@@M15@@

La section @@M16@@ est définie par un idéal isomorphe à @@M17@@ (4.2) ; on note @@M18@@ sa classe dans @@M19@@. On a alors la relation suivante pour tout @@M20@@ :
@@M21@@
@@M22@@ étant le morphisme structural @@M23@@. En effet, comme @@M24@@, on a @@M25@@, d'où @@M26@@; comme @@M27@@, on a la relation voulue :

La relation (4.3.3) s'écrit alors :
@@M28@@
